%\pdfoutput=1
\documentclass[10pt]{scrartcl}
\usepackage[T1]{fontenc}
\usepackage{lmodern}
\usepackage[ngerman]{babel}
\usepackage{tabto}
\usepackage{microtype}
%\usepackage{refcheck}
\setkomafont{author}{\small}
\sloppy
\emergencystretch=3em


\begin{document}
\title{\LARGE Vorhersage von Kundenentscheidungen: Ein Vergleich von sieben Klassifikatoren auf drei tabellarischen Datensätzen}
\author{
	Quirin Lex (xxx)\\
	Xuan Truong Nguyen (xxx)\\
	Johannes Ramser (xxx)\\
	Carlos García (xxx)
}
\date{}
\maketitle

\tableofcontents
\newpage

\begin{abstract}
\noindent
\textbf{Abstrakt:}
In dieser Arbeit handelt es sich um die binäre Klassifikation von Kundenentscheidungen auf tabellarischen Geschäftsdaten, nämlich die Vorhersage von Kündigung, Vertragsabschluss und Kaufabsicht. Die zentrale Herausforderung besteht darin, dass solche Datensätze typischerweise stark unbalancierte Klassen haben, was konventionelle Evaluationsmetriken wie Accuracy irreführend macht und die Modellwahl 
etwas erschwert. Der gewählte Ansatz umfasst einen systematischen Vergleich von sieben Klassifikationsverfahren — Logistic Regression, k-Nearest Neighbors, Naive Bayes, Support Vector Machine, Multi-Layer Perceptron, Random Forest und Gradient Boosting — auf drei öffentlich verfügbaren Datensätzen. Alle Modelle werden in einer leak-freien Preprocessing-Pipeline mit geschachtelter 
Kreuzvalidierung evaluiert, wobei F1-Score und ROC-AUC als Hauptmetriken dienen. Hier wurde es anhand der Ergebnisse gezeigt, dass ensemble-basierte Verfahren durchgehend die höchste Vorhersagegüte erzielen aber gleichzeitig eine erhöhte Anfälligkeit für Overfitting aufweisen. Eine ergänzende Modellanalyse mittels Permutation Importance, SHAP-Werten und Lernkurven verdeutlicht darüber hinaus, dass 
einzelne dominante Merkmale sowie starke Feature-Korrelationen die Interpretierbarkeit der Modelle erheblich einschränken können.
\end{abstract}


\section{Einleitung}
Heutzutage verfügen Unternehmen ständig über umfangreiche tabellarische Kundendaten, die von Vertragslaufzeiten und Nutzungsverhalten bis hin zu demografischen Merkmalen und makroökonomischen Indikatoren reichen. Dies führt dazu, dass die systematische Analyse und Auswertung solcher Daten erforderlich sind. Hierbei bieten die maschinellen Lernverfahren die Möglichkeit, Kundenentscheidungen frühzeitig 
vorherzusagen und daher lässt sich die Grundlage für gezielte Maßnahmen in Marketing, Kundenbindung und Ressourcenplanung bilden.

Die vorliegende Arbeit befasst sich mit der binären Klassifikation von Kundenentscheidungen auf tabellarischen Geschäftsdaten. Im Mittelpunkt steht die Frage, welche Klassifikationsverfahren unter praxisnahen Bedingungen, insbesondere bei stark unbalancierten Klassen, korrelierten Merkmalen und begrenztem Datenumfang, zuverlässige Vorhersagen liefern. Dabei geht es nicht nur um die reine Vorhersagegüte, 
sondern auch um die Interpretierbarkeit der Modelle und die Robustheit der Evaluationsergebnisse.

Die Relevanz dieser Fragestellung ergibt sich aus zwei Faktoren. Erstens ist die Wahl des Klassifikationsverfahrens in der Praxis häufig nicht trivial, d.h. unterschiedliche Modelle reagieren unterschiedlich auf Klassenverteilung und die Dimensionalität der Daten, sodass ein einzelner Datensatz oder eine einzelne Metrik kein vollständiges Bild liefert. Zweitens können fehlerhafte Evaluationsstrategien 
wie Data Leakage durch Preprocessing oder die Verwendung von Accuracy bei unbalancierten Klassen zu verzerrten Ergebnissen führen, die in der Praxis falsche Entscheidungen nach sich ziehen. Ein systematischer Vergleich unterschiedlicher Verfahren auf mehreren Datensätzen adressiert hierbei beide Probleme.

Vergleichende Studien zu Klassifikationsverfahren auf tabellarischen Geschäftsdaten sind in der Literatur breit vertreten. Arbeiten wie die von Fernández-Delgado et al.\cite{FernandezDelgado2014} untersuchen eine Vielzahl von Klassifikatoren auf zahlreichen Datensätzen und liefern wertvolle Orientierung bei der Modellwahl. Währenddessen stimmen datensatzspezifische Studien z. B. zur Vorhersage von Kundenabwanderung\cite{Ahmad2019} 
oder Kaufverhalten\cite{Sakar2019}, einzelne Verfahren gezielt auf bestimmte Anwendungskontexte ab. Die vorliegende Arbeit verfolgt einen ergänzenden Ansatz, indem sie einen einheitlichen methodischen Rahmen über mehrere Datensätze und Modelle hinweg anlegt. Daneben werden Aspekte wie Interpretierbarkeit, Generalisierungsverhalten und den Einfluss von Klassenungleichgewicht auch systematisch berücksichtigt.

Diese Arbeit leistet einen Beitrag, indem sie sieben Klassifikationsverfahren: Logistic Regression, k-Nearest Neighbors, Naive Bayes, Support Vector Machine, Multi-Layer Perceptron, Random Forest und Gradient Boosting, auf drei öffentlich verfügbaren Datensätzen (Telco Customer Churn, Bank Marketing, Online Shoppers Intention) vergleicht. Alle Modelle werden in einer einheitlichen, leak-freien Preprocessing-Pipeline 
mit geschachtelter Kreuzvalidierung evaluiert, wobei F1-Score und ROC-AUC als Hauptmetriken dienen. Über den reinen Modellvergleich hinaus umfasst die Analyse Lernkurven, Hyperparameter-Sensitivitätsanalysen, etc. , um nicht nur die Vorhersagegüte, sondern auch die Interpretierbarkeit und das Generalisierungsverhalten der Modelle zu bewerten.

\section{Verwandte Arbeiten}

\textbf{Gliederung (ca. 1 Seite):}
\begin{itemize}
	\item Literatur zu Churn-Prediction in Telekommunikation (klassische Verfahren: Logistic Regression, Decision Trees, Ensembles) — Einordnung des Telco-Datensatzes.
	\item Literatur zu Bank-Marketing/Term-Deposit-Vorhersage (UCI Bank-Marketing-Datensatz ist Standard-Benchmark) — welche Modelle/Resultate dort berichtet werden, insbesondere Rolle von \emph{duration} und makroökonomischen Indikatoren.
	\item Literatur zu Kaufabsicht/Online-Shopper-Intention-Vorhersage (Session-/Clickstream-basierte Merkmale, PageValues als bekannt starker Prädiktor).
	\item Allgemeine Methodik-Literatur: Umgang mit Klassenungleichgewicht (class\_weight, Resampling), Wahl von F1/ROC-AUC statt Accuracy, nested Cross-Validation zur unverzerrten Modellbewertung, SHAP/Permutation Importance als Standardverfahren der Modellinterpretation.
	\item Abgrenzung zur eigenen Arbeit: Vergleich über mehrere Domänen mit identischer Pipeline statt domänenspezifischer Einzelstudien; kritische Diskussion von Kollinearität (z.\,B. emp.var.rate vs. euribor3m, r=0.97) bei Feature-Importance-Methoden.
\end{itemize}

\section{Methoden}

\textbf{Gliederung (ca. 2--3 Seiten). Ziel: Studierende im selben Semester sollen die Kernidee nachvollziehen können.}

\subsection{Datenvorverarbeitung}
\label{sec:preprocessing}

Die drei Datensätze liegen in heterogenen Rohformaten mit gemischten Datentypen, domänenspezifischen Platzhaltern und unterschiedlichen Trennzeichen vor. Um eine reproduzierbare Vergleichsbasis zu schaffen, durchlaufen alle Datensätze eine identische Vorverarbeitungspipeline, die mit der Entfernung exakter Duplikate beginnt (22~Zeilen in Telco, 12 in Bank, 125 in Ecom). Spaltennamen werden vereinheitlicht, indem Leerzeichen und Schrägstriche durch Unterstriche ersetzt und Mehrfachunterstriche kollabiert werden; führende und nachfolgende Leerzeichen in String-Zellen werden bereinigt.

Nach der strukturellen Bereinigung erfolgt die Behandlung fehlender Werte. Die generischen Platzhalter \texttt{unknown}, \texttt{nonexistent}, \texttt{?} sowie leere Strings werden datensatzübergreifend durch \texttt{NaN} ersetzt. Der Wert \texttt{999} wird spaltenspezifisch ersetzt, da er kontextabhängig einen fehlenden Wert (Bank: 39.673\,Einträge in \texttt{pdays}) oder einen gültigen Messwert codiert. Die Spalte \texttt{default} im Bank-Datensatz enthält 8.597~\texttt{unknown}-Einträge.

Objektspalten, deren nicht-leere Werte zu mindestens 95\,\% in Zahlen umwandelbar sind, werden in einen numerischen Typ konvertiert; die Zielvariable ist ausgeschlossen. Verbleibende Fehlwerte werden imputiert: numerische Spalten erhalten den Median, kategoriale den Modus. Die Zielvariable wird abschließend über explizite Mapping-Tabellen binär kodiert (\texttt{0}/\texttt{1}).

Die Transformation für das Modelltraining erfolgt nicht global, sondern ausschließlich innerhalb von \texttt{Pipeline}-Objekten der Evaluierungsskripte. Numerische Merkmale werden mit \texttt{StandardScaler} standardisiert, kategoriale per \texttt{OneHotEncoder} (\texttt{handle\_unknown='ignore'}) One-Hot-kodiert; ein \texttt{ColumnTransformer} bettet beide als Preprocessing-Schritt in jede Modell-Pipeline ein. \texttt{fit()} wird nur auf den Trainingsfolds aufgerufen; Testdaten werden ausschließlich transformiert. Diese Trennung verhindert Data Leakage, das bei globaler Skalierung oder Kodierung vor der Aufteilung in Trainings- und Testdaten zu optimistisch verzerrten Evaluationsergebnissen führen würde.

Der Zufallsgenerator wird durchgängig mit \texttt{random\_state=42} fixiert.

\subsection{Baseline-Klassifikatoren (Task 2)}
\begin{itemize}
	\item Kurze formale Beschreibung von Logistic Regression (Sigmoid + Cross-Entropy-Loss).
	\item k-Nearest Neighbors (Distanzmaß, Mehrheitsentscheid, Wahl von k=5).
	\item Naive Bayes (Unabhängigkeitsannahme zwischen Features; Limitation bei korrelierten Bank-Features).
\end{itemize}

\subsection{Fortgeschrittene Modelle (Task 3)}
\label{sec:advanced-models}
\begin{itemize}
	\item Support Vector Machine: RBF-Kernel, Hyperparameter $C$ via GridSearchCV (Grid: $C \in \{0.01, 0.1, 1, 10\}$).
	\item Multi-Layer Perceptron: Architektur-Suche über \texttt{hidden\_layer\_sizes} $\in \{(64,), (128,), (64,32), (128,64)\}$ und Regularisierung $\alpha$; Early Stopping.
	\item Random Forest: Bagging von Entscheidungsbäumen, Hyperparameter \texttt{n\_estimators}, \texttt{max\_depth}, \texttt{max\_features}, \texttt{class\_weight=\textquotesingle balanced\textquotesingle} zur Behandlung des Klassenungleichgewichts.
	\item Gradient Boosting: konstruiert ein additives Ensemble flacher Entscheidungsbäume,
	die --- anders als beim Bagging des Random Forest --- sequentiell trainiert
	werden: Jeder Baum approximiert den negativen Gradienten der logistischen
	Verlustfunktion und korrigiert so die Fehler des bisherigen Modells, wodurch
	Boosting primär den Bias statt der Varianz reduziert \cite{Friedman2001}. Wir
	verwenden die histogramm-basierte, auf großen Datensätzen effiziente
	Implementierung \texttt{HistGradientBoostingClassifier} \cite{Pedregosa2011} in
	derselben leak-freien Pipeline wie die übrigen Modelle. Per \texttt{GridSearchCV}
	optimiert werden \texttt{learning\_rate} $\in \{0.01, 0.05, 0.1\}$,
	\texttt{max\_iter} $\in \{100, 200, 300\}$ und \texttt{max\_depth}
	$\in \{5, 10, 15, None\}$, wobei Lernrate und Iterationszahl den zentralen
	Regularisierungs-Kompromiss bilden.
	\item \textbf{Wichtig:} Hier nur die Kernidee jedes Verfahrens kurz erläutern (kein vollständiges Tutorial) — Fokus auf das tatsächlich beste Modell (Random Forest) etwas ausführlicher beschreiben, da Aufgabe explizit das \emph{beste} Modell in ausreichendem Detail verlangt.
\end{itemize}

\subsection{Modellauswahl und -bewertung}
\label{sec:model-selection}

Zur Bewertung der in Abschnitt~3.2 und~3.3 eingeführten Klassifikatoren wird jeder Datensatz in einen Trainings- (80\,\%) und einen Testanteil (20\,\%) aufgeteilt. Für Telco und Ecom erfolgt die Aufteilung stratifiziert über \texttt{train\_test\_split}, für Bank temporal an einem fixen Index, da die zeitlich geordneten makroökonomischen Merkmale einen Look-Ahead-Bias riskieren. Die Generalisierungsleistung aller sieben Modelle wird einheitlich mittels 5-facher Kreuzvalidierung auf dem Trainingsanteil geschätzt: \texttt{StratifiedKFold} ($k=5$) für Telco und Ecom erhält die Klassenverteilung pro Fold, \texttt{TimeSeriesSplit} ($k=5$) für Bank respektiert die zeitliche Ordnung. Für die fortgeschrittenen Modelle, deren Hyperparameter optimiert werden, kommt eine geschachtelte Kreuzvalidierung zum Einsatz: Ein innerer \texttt{GridSearchCV}-Loop bewertet Kandidaten anhand von \texttt{roc\_auc}, während ein äußerer Loop eine unverzerrte Schätzung der Generalisierungsfähigkeit liefert --- die äußeren Testdaten werden weder für Training noch für Tuning verwendet.

Als Evaluierungsmetriken dienen F1-Score und ROC-AUC als primäre, Accuracy als sekundäre Kennzahl. Der F1-Score als harmonisches Mittel aus Precision und Recall balanciert beide Größen und ist bei unbalancierten Klassen aussagekräftiger als Accuracy, die bereits ein trivialer Majority-Klassifikator hoch erscheinen ließe --- im Bank-Datensatz beträgt der Minoritätsanteil lediglich 11,3\,\%. Der ROC-AUC misst die Trennfähigkeit über verschiedene Entscheidungsschwellen hinweg und ist damit schwellenunabhängig. \texttt{random\_state=42} wird in sämtlichen datenaufteilenden und modelltrainierenden Operationen fixiert, um Reproduzierbarkeit zu gewährleisten.

\subsection{Modellanalyse-Framework (Task 4)}
\begin{itemize}
	\item \textbf{Motivation:} Reine Vorhersagegüte reicht für das Verständnis des Modells nicht aus --- das beste Modell (Random Forest / Bank Marketing) wird auf sechs Dimensionen vertieft analysiert, um Interpretierbarkeit, Generalisierungsverhalten und Fehlermuster zu bewerten.
	\item \texttt{model\_analysis.py}: 16 modell-agnostische Funktionen (8 Berechnung + 8 Plot) für das beste Modell, ausgeführt via \texttt{run\_analysis.py}.
	\item Kurze methodische Erläuterung der Analyse-Verfahren:
	\begin{itemize}
		\item \textbf{Permutation Importance} (Breiman, 2001): Feature-Werte werden zufällig permutiert; die Verschlechterung der Metrik (F1) misst die Wichtigkeit. Globales, modell-agnostisches Maß. Output: Balkendiagramm der Top-15 Features.
		\item \textbf{SHAP --- SHapley Additive exPlanations} (Lundberg \& Lee, 2017): Shapley-Werte aus der kooperativen Spieltheorie. Jedes Feature erhält einen additiven Beitrag zu jeder einzelnen Vorhersage (Summe aller SHAP-Werte + Baseline = Modellvorhersage). \texttt{TreeExplainer} für effiziente Berechnung bei baumbasierten Modellen. Output: Summary-Plot (Beeswarm, globale Wichtigkeit + Richtung) und Dependence-Plots (Top-5 Features). \textbf{Methodische Limitation:} Bei Multikollinearität (z.\,B. \texttt{emp.var.rate} vs. \texttt{euribor3m}, $r=0{,}97$) teilen sich korrelierte Features die SHAP-Werte --- die Importance ist konservativ und nicht kausal interpretierbar.
		\item \textbf{Lernkurven:} Trainings- vs. Validierungs-F1 über wachsende Trainingsgröße zur Diagnose von Bias vs. Varianz (Overfitting-Signal).
		\item \textbf{Hyperparameter-Sensitivitätsanalyse:} CV-Score als Funktion einzelner Hyperparameter (\texttt{n\_estimators}, \texttt{max\_depth}) zur Bewertung der Modell-Robustheit.
		\item \textbf{Konfusionsmatrix-Analyse} (absolut + normiert) und \textbf{Misclassification-Pattern-Analyse} (Feature-Mittelwerte: korrekt vs. falsch klassifiziert) zur Identifikation systematischer Fehlermuster.
	\end{itemize}
\end{itemize}

\section{Experimente}

\textbf{Gliederung (ca. 2--3 Seiten). Ziel der Experimente in 1--2 Sätzen formulieren:} Vergleich der Vorhersagegüte von sieben Klassifikatoren über drei Kundendatensätze hinweg unter einheitlicher, leak-freier Methodik, gefolgt von vertiefter Analyse des besten Modells (Random Forest / Bank Marketing) zur Erklärung seines Verhaltens.

\subsection{Daten}
\label{sec:data}

\begin{table}[ht]
\centering
\small
\caption{Übersicht der drei verwendeten Datensätze. Alle weisen ein Klassenungleichgewicht auf; Bank Marketing hat die geringste Minorität (11,3\,\%).}
\label{tab:datasets}
\begin{tabular}{lcccccc}
\hline
Datensatz & Quelle & Zeilen & Features & Zielklasse & Minorität \\
\hline
Telco Customer Churn & Kaggle/IBM & 7.043 & 20 & Churn & 26,5\,\% \\
Bank Marketing & UCI (Moro et al., 2014) & 41.188 & 21 & Term Deposit & 11,3\,\% \\
Online Shoppers Intention & UCI & 12.330 & 18 & Revenue & 15,5\,\% \\
\hline
\end{tabular}
\end{table}

\begin{itemize}
	\item Tabelle der drei Datensätze (siehe Tabelle~\ref{tab:datasets}).
	\item \textbf{Telco Customer Churn} (Kaggle/IBM, 7.043 Zeilen, 20 Features): Zielklasse \texttt{Churn}, Minorität 26,5\,\%. Fehlende Werte: \texttt{TotalCharges} (11 Zeilen). Duplikate: 22.
	\item \textbf{Bank Marketing} (UCI, Moro et al., 2014, 41.188 Zeilen, 21 Features): Zielklasse \texttt{Term Deposit}, Minorität 11,3\,\%. \textbf{Platzhalter-Problematik:} \texttt{job} (330$\times$ \texttt{unknown}), \texttt{education} (1.731$\times$), \texttt{default} (8.597$\times$), \texttt{pdays} (39.673$\times$ \texttt{999}). \textbf{Kollinearität:} \texttt{emp.var.rate} vs. \texttt{euribor3m} ($r=0{,}97$) --- relevant für die Interpretation der Feature-Importance in Abschnitt~\ref{sec:detailed-analysis}. Duplikate: 12.
	\item \textbf{Online Shoppers Intention} (UCI, 12.330 Zeilen, 18 Features): Zielklasse \texttt{Revenue}, Minorität 15,5\,\%. Platzhalter: \texttt{ProductRelated\_Duration} (3$\times$ \texttt{999}). Duplikate: 125.
	\item Metrik-Wahl: F1-Score und ROC-AUC als primäre Metriken statt Accuracy aufgrund des durchgängigen Klassenungleichgewichts in allen drei Datensätzen (vgl. Abschnitt~\ref{sec:model-selection}).
\end{itemize}

\subsection{Versuchsaufbau}
\begin{itemize}
	\item Algorithmen: Logistic Regression, k-NN (k=5), Naive Bayes (Baselines); SVM (RBF), MLP, Random Forest, Gradient Boosting (fortgeschritten).
	\item Metriken: Accuracy (zur Vollständigkeit), F1-Score und ROC-AUC (primär).
	\item Hyperparameter-Optimierung: GridSearchCV pro Modell (Grids siehe Abschnitt~\ref{sec:advanced-models}), Scoring \texttt{roc\_auc} im inneren CV-Loop.
	\item Modellauswahl/-bewertung: Nested Cross-Validation für fortgeschrittene Modelle; \texttt{StratifiedKFold} (Telco/Ecom) bzw. \texttt{TimeSeriesSplit} (Bank); finale Testresultate auf separatem Hold-out-Test-Set.
	\item Reproduzierbarkeit: fixer Random State 42, Pipelines mit \texttt{ColumnTransformer} dokumentiert in \texttt{src/evaluation.py}, \texttt{src/gradient\_boosting.py}, \texttt{src/run\_analysis.py}.
\end{itemize}

\subsection{Ergebnisse und Diskussion}

\paragraph{Hauptergebnisse — mit Tabellen/Abbildungen zu belegen.}
\begin{itemize}
	\item Tabelle: Baseline-CV-Ergebnisse (Accuracy/F1/ROC-AUC) für LogReg, k-NN, Naive Bayes auf allen 3 Datensätzen (Quelle: \texttt{reports/results/baseline\_cv\_results.csv}). Kernaussage: Logistic Regression dominiert die Baselines auf allen drei Datensätzen; Naive Bayes schwächste Baseline auf Bank (Unabhängigkeitsannahme verletzt durch korrelierte Features).
	\item Tabelle: Vergleich aller 7 Modelle nach ROC-AUC je Datensatz (Tab.~\ref{tab:overall}; Quellen: \texttt{baseline\_cv\_results.csv}, \texttt{gradient\_boosting\_results.csv}, \texttt{notebooks/Vergleich.ipynb}, \texttt{doc.md} Abschnitt 4).
	\item Abbildung: Balkendiagramm Validation-Accuracy aller Klassifikatoren je Datensatz (aus \texttt{Vergleich.ipynb}).
	\item Kernbefund: Random Forest auf Bank Marketing liefert höchste ROC-AUC (0.943--0.945, nested CV; Parameter: n\_estimators=200, max\_depth=20, max\_features='sqrt', class\_weight='balanced'); knapp gefolgt von MLP (0.941) und Gradient Boosting (0.916, Test-ROC-AUC).
	\item Gradient Boosting: erzielt über alle drei Datensätze hohe ROC-AUC-Werte (0.85--0.95, Tab.~\ref{tab:gb-results}) und übertrifft durchgängig die Baseline-Klassifikatoren. Auf Bank Marketing liegt die nested-CV-Schätzung (0.945) über der Test-ROC-AUC (0.914); diese Lücke erklärt sich durch den temporalen Split (Concept Drift bei den ökonomischen Indikatoren), nicht durch Overfitting. Die optimalen Hyperparameter für Telco und E-Commerce (niedrige Lernrate 0.01 bei 300 Iterationen) illustrieren den Regularisierungs-Kompromiss zwischen Lernrate und Iterationszahl.
\end{itemize}

\begin{table}[ht]
\centering
\small
\caption{Modellvergleich (ROC-AUC) über alle drei Datensätze. Baselines aus 5-fold CV, fortgeschrittene Modelle aus nested CV bzw. Hold-out-Test. \textit{tbd}: SVM-Werte (ROC-AUC/F1) von Florian noch nachzurechnen --- bislang nur Accuracy vorhanden. Bank-Spalte: MLP/RF aus Random-, GB aus Temporal-Split.}
\label{tab:overall}
\begin{tabular}{lccc}
\hline
Modell & Telco & Bank & E-Commerce \\
\hline
Logistic Regression & 0.846 & 0.928 & 0.890 \\
k-NN ($k=5$)        & 0.783 & 0.838 & 0.786 \\
Naive Bayes         & 0.821 & 0.819 & 0.812 \\
SVM (RBF)           & \textit{tbd} & \textit{tbd} & \textit{tbd} \\
MLP                 & 0.844 & 0.941 & 0.916 \\
Random Forest       & 0.838 & 0.943 & 0.924 \\
Gradient Boosting   & 0.848 & 0.945 & 0.932 \\
\hline
\end{tabular}
\end{table}

\begin{table}[ht]
\centering
\small
\caption{Gradient-Boosting-Ergebnisse je Datensatz (nested CV bzw. Hold-out-Test).
Beste Parameter (\texttt{learning\_rate}/\texttt{max\_depth}/\texttt{max\_iter}):
Telco 0.01/5/300, Bank 0.05/None/100, E-Commerce 0.01/None/300.}
\label{tab:gb-results}
\begin{tabular}{lcccc}
\hline
Datensatz & Split & Nested-CV ROC-AUC & Test ROC-AUC & Test F1 \\
\hline
Telco & random   & $0.848 \pm 0.010$ & 0.845 & 0.581 \\
Bank  & temporal & $0.945 \pm 0.014$ & 0.914 & 0.621 \\
Ecom  & random   & $0.932 \pm 0.005$ & 0.935 & 0.658 \\
\hline
\end{tabular}
\end{table}

\paragraph{Detaillierte Analyse des besten Modells (Random Forest / Bank Marketing) --- Modellanalyse Task 4.}
\label{sec:detailed-analysis}

Das beste Modell (Random Forest auf Bank Marketing, ROC-AUC 0,943) wird auf sechs Dimensionen vertieft analysiert. Ziel ist nicht nur die Bewertung der Vorhersagegüte, sondern das Verständnis von Feature-Wichtigkeit, Generalisierungsverhalten und Fehlermustern.

\subparagraph{Modellleistung und Permutation Importance.}
\begin{itemize}
	\item \textbf{Testergebnisse des finalen Modells} (Quelle: \texttt{reports/results/analysis\_results.csv}): ROC-AUC 0,912, F1 0,576, Precision 0,664, Recall 0,508, Specificity 0,949. Interpretation: ROC-AUC von 0,91 zeigt gute Trennfähigkeit; der moderate F1 spiegelt das Klassenungleichgewicht wider. Hohe Specificity bei moderatem Recall $\to$ das Modell ist konservativ bei der Vorhersage der positiven Klasse.
	\item \textbf{Permutation Importance:} \emph{duration} dominantes Feature (Importance 0,232), gefolgt von \texttt{emp.var.rate}, \texttt{euribor3m}, \texttt{nr.employed}, \texttt{poutcome}. Konsistent mit EDA: \texttt{duration} korreliert mit der Zielvariable ($r=0{,}405$). \textbf{Wichtig:} \texttt{duration} (Gesprächsdauer) ist nur post-hoc verfügbar --- für prospektive Vorhersagen vor einem Kundenkontakt nicht nutzbar, aber aufschlussreich für das Verständnis des Entscheidungsprozesses.
\end{itemize}

\subparagraph{SHAP-Analyse.}
\begin{itemize}
	\item \textbf{Summary-Plot:} SHAP bestätigt die globale Wichtigkeit von \texttt{duration}, liefert aber zusätzlich die Richtung des Einflusses (positive/negative SHAP-Werte) und lokale Erklärungen pro Instanz. Längere \texttt{duration} $\to$ höhere Wahrscheinlichkeit für Term Deposit.
	\item \textbf{Dependence-Plots:} Zeigen den Zusammenhang zwischen Feature-Wert und SHAP-Wert für die Top-5 Features (\texttt{duration}, \texttt{emp.var.rate}, \texttt{euribor3m}, \texttt{nr.employed}, \texttt{poutcome}).
	\item \textbf{Kollinearität als methodische Limitation:} \texttt{emp.var.rate} und \texttt{euribor3m} sind hoch korreliert ($r=0{,}97$, vgl. Abschnitt~\ref{sec:data}). Bei SHAP und Permutation Importance werden beide als wichtig ausgewiesen --- aber die Importance teilt sich zwischen kollinearen Features auf. Die SHAP-Werte sind deskriptiv, nicht kausal zu interpretieren. Dies ist eine bekannte Einschränkung von Shapley-basierten Verfahren unter Multikollinearität, kein Fehler des Modells.
\end{itemize}

\subparagraph{Lernkurven und Hyperparameter-Sensitivität.}
\begin{itemize}
	\item \textbf{Lernkurven:} Bias-Variance-Gap von 0,64 (F1-Differenz zwischen Trainings- und Validierungsleistung). Trainings-Accuracy $\approx$1,00 vs. Validierungs-Accuracy $\approx$0,91 $\to$ das Modell zeigt Anzeichen von Overfitting, was bei Random Forests auf Datensätzen dieser Größe zu erwarten ist. Die hohe Test-ROC-AUC (0,91) belegt, dass das Modell dennoch gut generalisiert. Weitere Daten oder stärkere Regularisierung könnten die Generalisierung verbessern.
	\item \textbf{Hyperparameter-Sensitivität:} \texttt{n\_estimators}: Plateau-Effekt ab $\approx$200 Bäume --- weitere Bäume bringen kaum Verbesserung. \texttt{max\_depth}: Tiefe $>$10 bringt marginalen Gewinn, erhöht aber das Overfitting-Risiko. Das Modell kann ohne signifikanten Performance-Verlust vereinfacht werden.
\end{itemize}

\subparagraph{Fehleranalyse.}
\begin{itemize}
	\item \textbf{Konfusionsmatrix:} Specificity 0,949, Recall 0,508. Das Modell findet etwa die Hälfte der tatsächlichen Term Deposits (Recall), produziert aber wenige False Positives. Geschäftsrelevanz: False Positives = unnötige Anrufe, False Negatives = verpasste Abschlüsse --- je nach Kostenstruktur kann der Entscheidungsschwellenwert angepasst werden.
	\item \textbf{Misclassification-Pattern:} Fehlklassifizierte Instanzen unterscheiden sich systematisch bei \texttt{duration} und den ökonomischen Indikatoren $\to$ das Modell ist insbesondere an der Entscheidungsgrenze (Fälle mit mittleren \texttt{duration}-Werten) unsicher.
\end{itemize}

\paragraph{Vergleich mit Literatur.} Einordnung der erreichten ROC-AUC-Werte (Bank $\approx$0.94) im Vergleich zu publizierten Resultaten auf dem UCI-Bank-Marketing-Datensatz; Diskussion möglicher Abweichungen (Preprocessing-Unterschiede, Wahl von TimeSeriesSplit statt zufälliger Split).

\paragraph{Limitationen (aus bestehender Discussion in \texttt{model\_analysis.ipynb} übernehmen und ausbauen).}
\begin{itemize}
	\item \textbf{\texttt{duration} nur post-hoc verfügbar:} \texttt{duration} (Gesprächsdauer) ist das mit Abstand wichtigste Feature, aber erst nach dem Kundenkontakt bekannt. Für prospektive Vorhersagen vor einem Anruf ist das Feature nicht verfügbar --- dies schränkt die praktische Anwendbarkeit des Modells für Targeting-Zwecke ein. Ein Modell ohne \texttt{duration} wäre für die Praxis relevanter, würde aber voraussichtlich niedrigere Vorhersagegüte aufweisen. Das Feature ist dennoch aufschlussreich für das Verständnis der Entscheidungsmechanismen.
	\item \textbf{Kollinearität und Feature-Importance:} \texttt{emp.var.rate} und \texttt{euribor3m} sind hoch korreliert ($r=0{,}97$). SHAP und Permutation Importance weisen beide als wichtig aus, können aber Kausalität nicht von Korrelation trennen --- die Importance teilt sich zwischen kollinearen Features auf und ist daher konservativ.
	\item Modellanalyse nur für einen Datensatz/ein Modell (Bank/Random Forest) durchgeführt; Telco und Ecom nicht vertieft analysiert.
	\item SHAP nur auf 200 Testsamples berechnet (Rechenaufwand) $\to$ mögliche Instabilität der Dependence-Plots, insbesondere bei seltenen Feature-Wert-Bereichen.
	\item Klassenungleichgewicht trotz \texttt{class\_weight='balanced'} weiterhin Herausforderung; SMOTE/Threshold-Tuning/Cost-Sensitive Learning nicht erprobt.
	\item Sensitivitätsanalyse nur für \texttt{n\_estimators}/\texttt{max\_depth}; weitere Parameter (\texttt{min\_samples\_leaf}, \texttt{min\_samples\_split}, \texttt{max\_features}) nicht untersucht.
	\item Bank-Datensatz zeitlich begrenzt (2008--2010); wirtschaftliche Veränderungen (Concept Drift) können die Vorhersagekraft des Modells beeinflussen --- \texttt{TimeSeriesSplit} adressiert dies methodisch, kann aber zukünftige Drifts nicht antizipieren.
\end{itemize}

\paragraph{Breitere Implikationen.} Geschäftsrelevanz der Precision/Recall-Charakteristik (Kosten von False Positives [unnötige Anrufe] vs.\ False Negatives [verpasste Abschlüsse]); Hinweis auf mögliche Schwellenwert-Anpassung je Geschäftsszenario.

\section{Fazit}

\textbf{Gliederung (1--2 Absätze, keine neuen Ergebnisse):}
\begin{itemize}
	\item Zusammenfassung: Random Forest erzielt über alle verglichenen Modelle und Datensätze hinweg die beste Generalisierungsleistung (Bank Marketing, ROC-AUC $\approx$0.94); Logistic Regression als stärkste Baseline; Klassenungleichgewicht und Feature-Korrelationen als durchgängige Herausforderungen.
	\item Limitationen kurz wiederholen (Overfitting/Bias-Variance-Gap, eingeschränkter Analyseumfang auf einen Datensatz, fehlende Exploration von Resampling-/Cost-Sensitive-Verfahren).
	\item Offene Fragen/Ausblick: Modellanalyse auf Telco/Ecom ausweiten, Threshold-/Cost-Sensitive-Tuning für Geschäftsszenarien, Untersuchung weiterer Hyperparameter, Robustheit gegenüber Concept Drift (Bank-Datensatz) genauer untersuchen.
\end{itemize}

\section{Literaturverzeichnis — Hinweis}

Alle in der Arbeit zitierten Quellen sind im Literaturverzeichnis aufzuführen.

Externe Aussagen, Gedanken und Ideen, die nicht zum Allgemeinwissen gehören (z.\,B. die Normalverteilung), müssen durch eine Quellenangabe belegt werden.

Bei der ersten Erwähnung von Verfahren, Konzepten, Frameworks usw. ist der Name der Autorin/des Autors in der Zitation anzugeben. Für Standardverfahren kann ein Lehrbuch als Quelle zitiert werden.

Der Zitier- und Referenzierungsstil ist über die gesamte Arbeit hinweg konsistent zu halten (z.\,B. APA, Chicago oder ein anderer gängiger Stil).

\paragraph*{Geplante Referenzen (Platzhalter, noch zu ergänzen):}
\begin{itemize}
	\item UCI Bank Marketing Dataset (Moro, Cortez, Rita) — Originalquelle des Bank-Datensatzes.
	\item Originalquelle Online Shoppers Purchasing Intention Dataset (Sakar et al.).
	\item Telco Customer Churn Dataset (IBM/Kaggle).
	\item Breiman (2001) — Random Forests.
	\item Lundberg \& Lee (2017) — SHAP (A Unified Approach to Interpreting Model Predictions).
	\item Lehrbuch-Standardreferenzen zu Logistic Regression, SVM, MLP und Gradient Boosting, z.\,B. Hastie/Tibshirani/Friedman, \emph{Elements of Statistical Learning}, oder Bishop, \emph{Pattern Recognition and Machine Learning}.
	\item Ggf. weitere Literatur zu Churn-Prediction/Bank-Marketing-Klassifikation aus der Related-Work-Recherche.
\end{itemize}

\paragraph*{Anmerkung:}
Die in diesem Dokument verwendete \verb|thebibliography|-Umgebung eignet sich für eine kleine Anzahl von Referenzen. Für größere Projekte sollten \verb|.bib|-Dateien zur besseren Verwaltung der Referenzen verwendet werden.



\addcontentsline{toc}{section}{Literaturverzeichnis}
\begin{thebibliography}{}
\small

\bibitem{FernandezDelgado2014}
M.~Fern\'{a}ndez-Delgado, E.~Cernadas, S.~Barro, and D.~Amorim.
\newblock Do we need hundreds of classifiers to solve real world classification problems?
\newblock \emph{Journal of Machine Learning Research}, 15(1):3133--3181, 2014.

\bibitem{Ahmad2019}
A.~K.~Ahmad, A.~Jafar, and K.~Aljoumaa.
\newblock Customer churn prediction in telecom using machine learning in big data platform.
\newblock \emph{Journal of Big Data}, 6(1):28, 2019.

\bibitem{Sakar2019}
C.~O.~Sakar, S.~O.~Polat, M.~Katircioglu, and Y.~Kastro.
\newblock Real-time prediction of online shoppers' purchasing intention using multilayer perceptron and LSTM recurrent neural networks.
\newblock \emph{Neural Computing and Applications}, 31(10):6893--6908, 2019.

% \bibitem{Breiman2001}
% L.~Breiman.
% \newblock Random forests.
% \newblock \emph{Machine Learning}, 45(1):5--32, 2001.

% \bibitem{Friedman2001}
% J.~H.~Friedman.
% \newblock Greedy function approximation: A gradient boosting machine.
% \newblock \emph{The Annals of Statistics}, 29(5):1189--1232, 2001.

% \bibitem{Lundberg2017}
% S.~M.~Lundberg and S.-I.~Lee.
% \newblock A unified approach to interpreting model predictions.
% \newblock \emph{Advances in Neural Information Processing Systems}, 2017.

% \bibitem{Moro2014}
% S.~Moro, P.~Cortez, and P.~Rita.
% \newblock A data-driven approach to predict the success of bank telemarketing.
% \newblock \emph{Decision Support Systems}, 62:22--31, 2014.

% \bibitem{Pedregosa2011}
% F.~Pedregosa et al.
% \newblock Scikit-learn: Machine learning in Python.
% \newblock \emph{Journal of Machine Learning Research}, 12:2825--2830, 2011.

\end{thebibliography}

\end{document}
